\subsection{Actual actor quotients preserve second collage cohomology}

All complexes in this subsection have homological grading.  The notation
$D[j]$ raises degrees by $j$.  All modules and homology groups are over
$\mathbb Q$.

\begin{lemma}[Proved: connective comparison after an idempotent-preserving quotient]
\label{CP2:A-idempotent-quotient-bridge-connectivity}
Let $\rho:M\twoheadrightarrow H$ be a surjective homomorphism of finite
monoids, and let $e\in M$ be idempotent.  Suppose the induced map
\[
 \mathbb QM e\longrightarrow\mathbb QH\rho(e)
                                                               \tag{1}
\]
is an isomorphism of left $\mathbb QM$ modules, where the action on the
right is through $\rho$.
For every finite left-$N$, right-$M$ biset $S$, form the actual induced
left-$N$, right-$H$ set
\[
 U=S\otimes_M H.
\]
Then the actual bridge map
\[
 B(N,S,M)\longrightarrow B(N,U,H)
                                                               \tag{2}
\]
induces isomorphisms on rational homology in degrees zero and one,
and a surjection in degree two.
The actual corner map $BM\to BH$ has the same low-degree properties.
\end{lemma}
\begin{proof}
Put $R=\mathbb QM$, $B=\mathbb QH$, $T=\mathbb QN$, and write $e$
for the image idempotent as well.  Let $L=\mathbb Q[S]$ and
$L'=\mathbb Q[U]$.
Linearization preserves the defining set coequalizer, so
\[
 L'=L\otimes_R B.                                               \tag{3}
\]
Since $\rho$ is surjective, every state $[s,h]$ of $U$ is of the form
$[sm,1]$ with $\rho(m)=h$.  Thus $U$ is finite and is an actual quotient
of $S$.  Its commuting $N$ action and the map $S\to U$ are canonical.

Use the ordinary nonnegative two-sided bar complex to represent
\[
 F=L\otimes_R^{\mathbf L}B
\]
as a complex of left-$T$, right-$B$ bimodules.  Its canonical
augmentation to (3) is degreewise surjective and induces the identity
on degree-zero homology.  Let $D$ be its kernel complex.  Consequently
\[
 H_i(D)=0\quad(i\leq0),\qquad
 D\longrightarrow F\longrightarrow L'
                                                               \tag{4}
\]
is the natural fiber triangle.  The usual good truncation replaces
$D_1$ by $\ker(D_1\to D_0)$ and deletes $D_0$; it is a natural
quasi-isomorphism from a complex in degrees at least one to $D$.
No splitting as a left-$T$ module is needed.

Multiplication by the idempotent on the right is exact.  By (1),
\[
 Fe\simeq L\otimes_R^{\mathbf L}Be
      \simeq L\otimes_R Re=Le
\]
is concentrated in degree zero, because $Re$ is left projective.
Ordinary tensor associativity and (3) give $L'e\cong Le$ as well.
The map $Fe\to L'e$ is this same natural isomorphism on degree-zero
homology.  Therefore
\[
 De\simeq0.                                                    \tag{5}
\]

The augmentation of the left projective $Be$ is surjective.  Write its
kernel as $V$, giving the short exact sequence
\[
 0\longrightarrow V\longrightarrow Be\longrightarrow\mathbb Q
          \longrightarrow0.                                  \tag{6}
\]
Derived tensoring (6) with $D$, and using
$D\otimes_B^{\mathbf L}Be\simeq De\simeq0$, gives
\[
 D\otimes_B^{\mathbf L}\mathbb Q
       \simeq (D\otimes_B^{\mathbf L}V)[1].                    \tag{7}
\]
The complex $D$ is in homological degrees at least one, and $V$ is a
module in degree zero.  Bounded-below free resolutions, or the
nonnegative two-sided bar, show that their derived tensor still has
no homology below degree one.  The right side of (7) therefore has
no homology below degree two.  Applying the nonnegative derived
$N$-coinvariants preserves that degree bound.

The natural derived associativity comparison identifies the chains of
the actual map (2) with
\[
 \mathbb Q\otimes_T^{\mathbf L}L\otimes_R^{\mathbf L}\mathbb Q
 \simeq
 \mathbb Q\otimes_T^{\mathbf L}F\otimes_B^{\mathbf L}\mathbb Q
 \longrightarrow
 \mathbb Q\otimes_T^{\mathbf L}L'\otimes_B^{\mathbf L}\mathbb Q.
                                                               \tag{8}
\]
Indeed the bar map is induced by the actual unit $s\mapsto[s,1]$
and by $\rho$ on every monoid arrow; its adjoint is precisely the
augmentation in (4).  Thus its fiber is
$\mathbb Q\otimes_T^{\mathbf L}D\otimes_B^{\mathbf L}\mathbb Q$,
which has no homology below degree two.  The homology exact sequence
gives the asserted isomorphisms and surjection.
Finally take $L=\mathbb Q$ with its trivial right-$R$ action and omit
the left $N$ bar.  Surjectivity of $\rho$ identifies
$\mathbb Q\otimes_R B$ with the trivial right-$B$ module, so the
same argument is exactly the actual corner comparison.
\end{proof}

\begin{proposition}[Proved: second cohomology of a collage under a band-preserving quotient]
\label{CP2:A-band-quotient-collage-H2}
Let $\rho:M\twoheadrightarrow H$ identify the proper rectangular
minimal ideals of the two finite monoids bijectively.  For every
finite left-$N$, right-$M$ biset $S$, the actual quotient functor
\[
 \mathcal C(N,S,M)\longrightarrow
       \mathcal C(N,S\otimes_M H,H)                            \tag{9}
\]
induces an isomorphism in rational second homology and cohomology.
In particular, if all degree-two products vanish on the target,
they vanish on the source.
\end{proposition}
\begin{proof}
Choose $e=e_{00}$ in the identified rectangular band $I\times J$.
The sets $Me$ and $He$ are its same column of $|I|$ elements.
Thus $Re\to Be$ is a left-$R$ isomorphism, and the preceding lemma
applies to the bridge and corner maps.

For completeness, the corner map is also an isomorphism in degree
two, rather than merely a surjection.  Let
\[
 V_M=\ker(\mathbb QI\to\mathbb Q),\qquad
 W_M=\ker(\mathbb QJ\to\mathbb Q).
\]
The two natural augmentation shifts of Lemma~\ref{A:band-Tor} identify
$H_2(BM;\mathbb Q)$ with the balanced quotient of $W_M\otimes V_M$.
The quotient monoid has the same two modules and exactly the same
row and column actions: these actions are read from multiplication
on the identified band.  Since $\rho$ is surjective, the balancing
relations on both sides are identical.  Naturality of the two
augmentation sequences therefore makes $BM\to BH$ an isomorphism
on $H_2$, and the preceding lemma already gives the needed $H_1$
isomorphism.

Use the natural two-block homology sequence.  Its middle portion is
\[
 H_2(E)\longrightarrow H_2(BM)\oplus H_2(BN)
 \longrightarrow H_2(B\mathcal C)
 \longrightarrow H_1(E)
 \longrightarrow H_1(BM)\oplus H_1(BN).
                                                               \tag{10}
\]
For (9), the map on the first term is surjective, the map on the
second is an isomorphism, and the maps on the last two terms are
isomorphisms.  Hence both the cokernel preceding $H_2(B\mathcal C)$
and its following kernel are isomorphic to their target versions.
The associated short exact sequences imply that the middle map is
an isomorphism as well.  Over $\mathbb Q$ the cohomological pullback
is its dual and is therefore an isomorphism.  It is induced by the
actual functor (9), so it preserves cup products.
If $S$ is empty, (9) is the disjoint union of the corner maps and
the same conclusion follows directly.
\end{proof}

\begin{theorem}[Proved: every two-by-two minimal band opposite $R_3$]
\label{CP2:A-two-by-two-band-right-zero-all-states}
Let $M$ be any finite monoid whose proper minimal ideal is a
$2\times2$ rectangular band.  Let $R_3$ act on the left of any finite
right-$M$ connecting set.  Every rational degree-two product on the
resulting collage is zero.  No bound on $|M|$, on the number of states,
or on the homology of the original collage is required.
\end{theorem}
\begin{proof}
Write the band as $\{e_{ij}:i,j\in\{0,1\}\}$.
For each $m\in M$ its two actors are the transformations
\[
 me_{ij}=e_{\sigma_m(i),j},\qquad
 e_{ij}m=e_{i,\tau_m(j)}.
\]
On the one-dimensional row and column augmentation spaces these
act by characters $\chi_I(m)$ and $\chi_J(m)$, respectively.
A transformation of a two-element set has augmentation scalar $1$
if it is the identity, $-1$ if it is the transposition, and $0$ if
it is constant.  The balanced description in the preceding proof
therefore says
\[
 H_2(BM;\mathbb Q)=
 \mathbb Q\big/\langle\chi_I(m)-\chi_J(m):m\in M\rangle.
                                                               \tag{11}
\]
If this group is zero, the other corner $BR_3$ is contractible, so
every degree-two class is relative to the two corners.  The actual
suspension quotient in the two-block model kills all their products.
Assume from now on that (11) is nonzero.  Then the two characters
agree for every element, and the group in (11) is one-dimensional.

Consider the actual actor homomorphism
\[
 \rho:M\longrightarrow H,\qquad
 m\longmapsto(\sigma_m,\tau_m),
\]
where actor pairs multiply by
$(\sigma,\tau)(\sigma',\tau')=(\sigma\sigma',\tau'\tau)$,
and $H$ is its image.  Character equality allows exactly six actor
pairs: the four pairs of constant transformations, the pair of
identities, and the pair of transpositions.  The four constant pairs
are present because they are the images of the four distinct band
elements; the identity pair is present because $\rho$ is unital.
Thus $H$ is either $P_5$, or the six-element monoid $P_6$ obtained by
adjoining a unit $g$ to the same band with
\[
 g^2=1,\qquad ge_{ij}=e_{1-i,j},\qquad e_{ij}g=e_{i,1-j}.
                                                               \tag{12}
\]
The actor map is bijective on the band, which remains the proper
minimal ideal in either image.  Proposition~\ref{CP2:A-band-quotient-collage-H2}
gives an actual quotient of the original collage onto a collage
with source $H$, inducing an isomorphism in degree-two cohomology.

Here is an explicit check of the target homological bound needed by
the existing graph theorem.  Put $B=\mathbb QH$ and
$A=B/\mathbb QK$.  For $P_5$ the algebra $A$ is $\mathbb Q$;
for (12) it is $\mathbb Q[C_2]$, which splits as
$\mathbb Q\oplus\mathbb Q$ using $(1+g)/2$ and $(1-g)/2$.
In either case $A$ is semisimple.
The band ideal has the left-projective factorization
\[
 \mathbb QK\cong Be\otimes_{\mathbb Q}eB.
\]
Consequently $0\to\mathbb QK\to B\to A\to0$ is a left-projective
resolution of $A$ of length one.  The row augmentation module $V$
is a projective $A$ module, so it has left-$B$ projective dimension
at most one.  The exact sequence
$0\to V\to Be\to\mathbb Q\to0$ then gives
\[
 \operatorname{pd}_B(\mathbb Q)\leq2.                           \tag{13}
\]
The trivial right $\mathbb Q R_3$ module has projective dimension
at most one, by its displayed right-zero resolution in
Corollary~\ref{CP2:A-baseline-right-zero-three-six}.
The target double-bar bridge consequently has rational homology
only through degree three, and its two corners only through degree
two.  The two-block sequence puts all target collage homology in
degrees at most four, for every finite induced set $U$.

We can therefore apply exactly
Corollary~\ref{CP2:A-baseline-right-zero-three-six}: it covers any
source monoid of order at most six, opposite $R_3$, with arbitrary
state count and bounded rational collage homology.  Both image monoids
satisfy its hypotheses by (12)--(13).  All target degree-two products
vanish.  Every second-degree class on the original collage is their
actual pullback, so all its degree-two products vanish as well.
If the connecting set is empty, the target is $BH\sqcup BR_3$;
(13) gives $H^4(BH;\mathbb Q)=0$, and the same pullback conclusion
handles that case directly.
\end{proof}
